\documentclass[9pt,a4paper]{article}

% ---- Paquetes ----
\usepackage[utf8]{inputenc}
\usepackage[T1]{fontenc}
\usepackage[spanish]{babel}
\usepackage{amsmath,amssymb,amsfonts}
\usepackage[margin=1.5cm]{geometry}
\usepackage{hyperref}
\usepackage{physics}

% Compactar espaciado
\setlength{\parskip}{1pt}
\setlength{\parindent}{0pt}
\usepackage{titlesec}
\titlespacing*{\section}{0pt}{5pt}{2pt}
\titlespacing*{\subsection}{0pt}{3pt}{1pt}

\title{\textbf{Separata de Estudio}\\[2pt]
\large Resolución de Sistemas de ED Lineales 2×2}
\author{Métodos: Eliminación, Operador Diferencial y Transformada de Laplace}
\date{}

\begin{document}
\maketitle

% ============================================================
\section*{Sistema general}
% ============================================================

Todo sistema lineal de dos ecuaciones con coeficientes constantes se escribe como:

\[
\boxed{
\begin{cases}
x'(t) = a_{11}\,x + a_{12}\,y \\[2pt]
y'(t) = a_{21}\,x + a_{22}\,y
\end{cases}
\qquad\text{con}\qquad
\begin{aligned}
x(0) &= x_0 \\
y(0) &= y_0
\end{aligned}
}
\]

\smallskip
Los tres métodos siguientes funcionan para \textbf{cualquier} sistema de esta forma. Solo cambian los coeficientes \(a_{ij}\) y las condiciones iniciales.


% ============================================================
\section{Método 1: Eliminación}
% ============================================================

\subsection{Paso 1 --- Despejar una variable}

De la primera ecuación, despejamos \(y\) en términos de \(x\) y \(x'\):

\[
y = \frac{x' - a_{11}\,x}{a_{12}}, \qquad a_{12} \neq 0
\]

Si \(a_{12} = 0\), se despeja de la segunda ecuación.

\subsection{Paso 2 --- Derivar}

Derivamos la expresión obtenida:

\[
y' = \frac{x'' - a_{11}\,x'}{a_{12}}
\]

\subsection{Paso 3 --- Sustituir en la otra ecuación}

Reemplazamos \(y\) e \(y'\) en la segunda ecuación del sistema:

\[
\frac{x'' - a_{11}\,x'}{a_{12}} = a_{21}\,x + a_{22}\left(\frac{x' - a_{11}\,x}{a_{12}}\right)
\]

\subsection{Paso 4 --- Obtener una EDO de segundo orden}

Multiplicamos por \(a_{12}\) y agrupamos términos:

\[
x'' - a_{11}\,x' = a_{12}\,a_{21}\,x + a_{22}\,x' - a_{22}\,a_{11}\,x
\]

\[
\boxed{x'' - (a_{11}+a_{22})\,x' + (a_{11}a_{22} - a_{12}a_{21})\,x = 0}
\]

Es decir:

\[
x'' - (\text{traza})\,x' + (\text{determinante})\,x = 0
\]

\subsection{Paso 5 --- Resolver la EDO homogénea}

Ecuación característica:

\[
r^{2} - (a_{11}+a_{22})\,r + (a_{11}a_{22} - a_{12}a_{21}) = 0
\]

Según el discriminante \(\Delta = (a_{11}+a_{22})^{2} - 4(a_{11}a_{22} - a_{12}a_{21})\):

\begin{itemize}
  \item[\textbf{Caso 1}] \(\Delta > 0\): dos raíces reales distintas \(r_1 \neq r_2\)
  \[
  x(t) = C_1 e^{r_1 t} + C_2 e^{r_2 t}
  \]

  \item[\textbf{Caso 2}] \(\Delta = 0\): raíz real doble \(r\)
  \[
  x(t) = C_1 e^{r t} + C_2 t e^{r t}
  \]

  \item[\textbf{Caso 3}] \(\Delta < 0\): raíces complejas \(r = \alpha \pm i\beta\)
  \[
  x(t) = e^{\alpha t}\bigl(C_1 \cos\beta t + C_2 \sin\beta t\bigr)
  \]
\end{itemize}

\subsection{Paso 6 --- Hallar la otra variable}

Volvemos a la relación del Paso 1:

\[
y = \frac{x' - a_{11}\,x}{a_{12}}
\]

Derivamos \(x(t)\) obtenida en el Paso 5, sustituimos y simplificamos.

\subsection{Paso 7 --- Solución final}

\[
\boxed{
\begin{aligned}
x(t) &= \text{solución del Paso 5} \\
y(t) &= \text{resultado del Paso 6}
\end{aligned}
}
\]

Si hay condiciones iniciales, se resuelve para \(C_1, C_2\).

% ============================================================
\section{Método 2: Operador Diferencial}
% ============================================================

\subsection{Paso 1 --- Escribir con el operador \(D = \frac{d}{dt}\)}

\[
\begin{aligned}
(D - a_{11})\,x - a_{12}\,y &= 0 \\
-a_{21}\,x + (D - a_{22})\,y &= 0
\end{aligned}
\]

\subsection{Paso 2 --- Formar la matriz del sistema}

\[
\begin{pmatrix}
D - a_{11} & -a_{12} \\
-a_{21}    & D - a_{22}
\end{pmatrix}
\begin{pmatrix} x \\ y \end{pmatrix}
=
\begin{pmatrix} 0 \\ 0 \end{pmatrix}
\]

\subsection{Paso 3 --- Calcular el determinante operacional}

\[
\Delta = (D - a_{11})(D - a_{22}) - (-a_{12})(-a_{21})
       = D^{2} - (a_{11}+a_{22})D + (a_{11}a_{22} - a_{12}a_{21})
\]

\subsection{Paso 4 --- Ecuación característica}

\[
r^{2} - (a_{11}+a_{22})\,r + (a_{11}a_{22} - a_{12}a_{21}) = 0
\]

(Idéntica a la de eliminación.)

\subsection{Paso 5 --- Resolver raíces}

Mismos 3 casos que en eliminación (Paso 5 del Método 1).

\subsection{Paso 6 --- Solución general para \(x(t)\)}

Según el caso de raíces:

\[
x(t) =
\begin{cases}
C_1 e^{r_1 t} + C_2 e^{r_2 t}, & \text{raíces reales distintas} \\[2pt]
C_1 e^{r t} + C_2 t e^{r t},   & \text{raíz real doble} \\[2pt]
e^{\alpha t}(C_1\cos\beta t + C_2\sin\beta t), & \text{raíces complejas}
\end{cases}
\]

\subsection{Paso 7 --- Sustituir en una ecuación original}

Usamos, por ejemplo, la primera ecuación:

\[
y = (D - a_{11})\,x = x' - a_{11}\,x
\]

Derivamos \(x(t)\) y calculamos \(y\).

\subsection{Paso 8 --- Relacionar constantes (si es necesario)}

En algunos casos al simplificar \(y\) aparecen las mismas constantes \(C_1, C_2\) ya combinadas. Se dejan expresadas.

\subsection{Paso 9 --- Solución final}

\[
\boxed{
\begin{aligned}
x(t) &= \text{del Paso 6} \\
y(t) &= \text{del Paso 7}
\end{aligned}
}
\]

% ============================================================
\section{Método 3: Transformada de Laplace}
% ============================================================

\subsection{Paso 1 --- Aplicar Laplace a cada ecuación}

Recordando \(\mathcal{L}\{x'\} = sX(s) - x(0)\):

\[
\begin{aligned}
sX - x_0 &= a_{11}X + a_{12}Y \\
sY - y_0 &= a_{21}X + a_{22}Y
\end{aligned}
\]

\subsection{Paso 2 --- Reordenar el sistema algebraico}

\[
\begin{cases}
(s - a_{11})\,X - a_{12}\,Y = x_0 \\[2pt]
-a_{21}\,X + (s - a_{22})\,Y = y_0
\end{cases}
\]

\subsection{Paso 3 --- Resolver para \(X(s)\) y \(Y(s)\)}

Se puede resolver por:
\begin{itemize}
  \item Sustitución (despejar una variable de una ecuación y reemplazar en la otra)
  \item Regla de Cramer:
  \[
  \Delta(s) = (s - a_{11})(s - a_{22}) - a_{12}a_{21}
  \]
  \[
  X(s) = \frac{
    \begin{vmatrix}
    x_0 & -a_{12} \\
    y_0 & s - a_{22}
    \end{vmatrix}
  }{\Delta(s)},
  \qquad
  Y(s) = \frac{
    \begin{vmatrix}
    s - a_{11} & x_0 \\
    -a_{21}    & y_0
    \end{vmatrix}
  }{\Delta(s)}
  \]
\end{itemize}

\subsection{Paso 4 --- Simplificar con fracciones parciales}

Se descompone cada expresión en fracciones simples según las raíces del polinomio \(\Delta(s)\):

\[
\frac{P(s)}{(s - r_1)(s - r_2)} = \frac{A}{s - r_1} + \frac{B}{s - r_2}
\]

\[
\frac{P(s)}{(s - r)^2} = \frac{A}{s - r} + \frac{B}{(s - r)^2}
\]

\[
\frac{P(s)}{(s - \alpha)^2 + \beta^2} = \frac{A(s - \alpha) + B\beta}{(s - \alpha)^2 + \beta^2}
\]

\subsection{Paso 5 --- Aplicar transformada inversa}

Usamos la tabla de transformadas inversas:

\[
\begin{aligned}
\mathcal{L}^{-1}\left\{\frac{1}{s - r}\right\} &= e^{r t} \\[2pt]
\mathcal{L}^{-1}\left\{\frac{1}{(s - r)^2}\right\} &= t e^{r t} \\[2pt]
\mathcal{L}^{-1}\left\{\frac{s - \alpha}{(s - \alpha)^2 + \beta^2}\right\} &= e^{\alpha t}\cos\beta t \\[2pt]
\mathcal{L}^{-1}\left\{\frac{\beta}{(s - \alpha)^2 + \beta^2}\right\} &= e^{\alpha t}\sin\beta t
\end{aligned}
\]

\subsection{Paso 6 --- Solución particular}

\[
\boxed{
\begin{aligned}
x(t) &= \mathcal{L}^{-1}\{X(s)\} \\
y(t) &= \mathcal{L}^{-1}\{Y(s)\}
\end{aligned}
}
\]

% ============================================================
\section*{Resumen: cómo elegir el método}
% ============================================================

\begin{tabular}{p{3cm} p{10cm}}
\textbf{Método} & \textbf{Cuándo usarlo} \\ \hline
Eliminación     & Cuando una variable se despeja fácilmente. Bueno para sistemas pequeños (2×2). No necesita condiciones iniciales. \\
Operador \(D\)  & Rápido cuando el determinante operacional se calcula sin esfuerzo. Muy sistemático, fácil de recordar. No necesita condiciones iniciales. \\
Laplace         & \textbf{Único} que da la solución particular directamente con las CI. Ideal si el examen \textbf{pide} Laplace. Útil cuando hay términos no homogéneos.
\end{tabular}

\medskip
\centerline{\textbf{¡Los tres métodos deben dar la misma solución!}}
\centerline{Úsalo como verificación: si discrepan, revisa las cuentas.}

\end{document}
